\documentclass[11pt]{beamer}

\usetheme{Madrid}
\usecolortheme{default}
\definecolor{unicaucaverde}{RGB}{0,90,60}
\setbeamercolor{structure}{fg=unicaucaverde}
\setbeamertemplate{navigation symbols}{}

\usepackage[utf8]{inputenc}
\usepackage[spanish]{babel}
\usepackage{amsmath,amssymb}
\usepackage{listings}
\usepackage{xcolor}

\lstdefinestyle{cstyle}{
  language=C,
  basicstyle=\ttfamily\scriptsize,
  keywordstyle=\color{unicaucaverde}\bfseries,
  commentstyle=\color{gray}\itshape,
  stringstyle=\color{orange!70!black},
  breaklines=true,
  showstringspaces=false,
  tabsize=2
}
\lstset{style=cstyle}

\title[Semana 1: Fundamentos de C]{Introducción a la programación y fundamentos de C}
\subtitle{Programación --- Semana 1 de 16}
\author{Universidad del Cauca}
\date{2026}

\begin{document}

\begin{frame}
  \titlepage
\end{frame}

\begin{frame}{Semana 1: dónde estamos}
  \begin{itemize}
    \item Unidad: Fundamentos de programación en C.
    \item Subtemas 1.1--1.4 --- RA1.1, RA1.2.
    \item 4 horas: proceso de compilación y primer programa, luego tipos/operadores/E-S.
  \end{itemize}
  \tableofcontents
\end{frame}

\section{Algoritmos y lenguajes}

\begin{frame}{¿Qué es programar?}
  \begin{block}{Algoritmo}
    Secuencia finita, ordenada y no ambigua de pasos que resuelve un problema.
  \end{block}
  \begin{block}{Programa}
    Traducción de un algoritmo a un lenguaje que el computador puede ejecutar.
  \end{block}
  \vfill
  \textbf{Compilado} (C): se traduce todo antes de ejecutar $\to$ ejecutable nativo.\\
  \textbf{Interpretado} (Python): se traduce y ejecuta línea a línea.
\end{frame}

\section{Compilación y estructura de un programa}

\begin{frame}{El proceso de compilación en C}
  \begin{enumerate}
    \item \textbf{Preprocesador} --- procesa \texttt{\#include}, \texttt{\#define}.
    \item \textbf{Compilador} --- código fuente $\to$ ensamblador.
    \item \textbf{Ensamblador} --- ensamblador $\to$ código objeto.
    \item \textbf{Enlazador} --- objeto + librerías $\to$ ejecutable.
  \end{enumerate}
  \vfill
  \texttt{gcc programa.c -o programa} hace las cuatro etapas de un solo paso.
\end{frame}

\begin{frame}[fragile]{Hola mundo}
\begin{lstlisting}
#include <stdio.h>

int main(void) {
    printf("Hola, mundo\n");
    return 0;
}
\end{lstlisting}
  \begin{itemize}
    \item \texttt{\#include}: directiva de preprocesador.
    \item \texttt{main}: punto de entrada del programa.
    \item Cada sentencia termina en \texttt{;}; el cuerpo va entre \texttt{\{ \}}.
  \end{itemize}
  \vfill
  \textbf{Taller (hora 2):} compilar, ejecutar, modificar, provocar/corregir errores.
\end{frame}

\section{Tipos de datos, variables y constantes}

\begin{frame}{Tipos de datos básicos}
  \begin{center}
  \begin{tabular}{lll}
    \textbf{Tipo} & \textbf{Representa} & \textbf{Literal} \\ \hline
    \texttt{int} & entero & \texttt{42} \\
    \texttt{float} & real simple & \texttt{3.14f} \\
    \texttt{double} & real doble & \texttt{3.14159265} \\
    \texttt{char} & carácter & \texttt{'A'} \\
  \end{tabular}
  \end{center}
  \vfill
  El tamaño exacto en bytes se consulta con \texttt{sizeof}.
\end{frame}

\begin{frame}[fragile]{Variables y constantes}
\begin{lstlisting}
int edad = 20;
const float PI = 3.1416f;
#define MAX_ESTUDIANTES 40
\end{lstlisting}
  \begin{itemize}
    \item Variable: declarar tipo + nombre, puede inicializarse.
    \item \texttt{const}: constante con tipo, verificada por el compilador.
    \item \texttt{\#define}: constante por sustitución de texto (preprocesador).
  \end{itemize}
\end{frame}

\section{Operadores y entrada/salida}

\begin{frame}{Familias de operadores}
  \begin{itemize}
    \item \textbf{Aritméticos:} \texttt{+ - * / \%}
    \item \textbf{Relacionales:} \texttt{== != < > <= >=}
    \item \textbf{Lógicos:} \texttt{\&\& || !}
    \item \textbf{Asignación:} \texttt{= += -= *= /=}
    \item \textbf{Incremento/decremento:} \texttt{++ --}
  \end{itemize}
  \begin{alertblock}{Cuidado}
    \texttt{/} entre dos \texttt{int} trunca: \texttt{7/2} vale \texttt{3}, no \texttt{3.5}.
  \end{alertblock}
\end{frame}

\begin{frame}[fragile]{\texttt{printf} y \texttt{scanf}}
\begin{lstlisting}
float radio, area;
printf("Radio: ");
scanf("%f", &radio);
area = 3.14159f * radio * radio;
printf("Area = %.2f\n", area);
\end{lstlisting}
  \begin{itemize}
    \item \texttt{\%d} entero, \texttt{\%f}/\texttt{\%lf} flotante, \texttt{\%c} carácter.
    \item \texttt{scanf} necesita \texttt{\&} (dirección) delante de cada variable.
  \end{itemize}
  \vfill
  \textbf{Taller (hora 4):} programa integrador con entrada de datos.
\end{frame}

\begin{frame}{Cierre de la semana}
  \begin{itemize}
    \item RA1.1: proceso de compilación, estructura mínima de un programa.
    \item RA1.2: tipos de datos, variables, constantes, operadores, E/S estándar.
    \item Control formativo al inicio de la Semana 2.
  \end{itemize}
  \vfill
  \centering \Large Semana 2: estructuras de control --- condicionales y ciclos.
\end{frame}

\end{document}
